% GENERATED FILE. Edit canonical lessons, then run npm run generate:books.
% canonical-source-hash: fnv1a64:d5048bdbd5a79ad2
% canonical-lessons: ZH-C295-zuichu, ZH-C295-chongxin, ZH-W295-xin-new, ZH-C295-dangtian, ZH-C295-shihou, ZH-C295-dangchang, ZH-R295-review-290-295

\chapter{Starting Over}
\label{ch:zh-starting-over}
\hlchaptermodality{295}

\begin{chapteropening}
\textbf{By the end of this chapter:} \emph{I can say how something was at first, that it began again, and that it happened that same day.}

Complete the last lesson of chapter 295: I can say how something was at first, that it began again, and that it happened that same day.
\end{chapteropening}

\section[zuìchū]{zuìchū — at the very start; initially}

\label{lesson:ZH-C295-zuichu}



\begin{quote}
Before the new one: say the Mandarin for at long last, then the Mandarin for and so it followed.
\end{quote}

\subsection*{You'll want to know: zuìchū}
\emph{zuìchū} — \textquotedblleft{}at the very start; initially\textquotedblright{}.

\begin{quote}
\emph{zuìchū wǒ bù xǐhuan zhèr} — At first I didn't like it here.
\end{quote}

Its characters come later; for now it is a word for the ear.

\begin{cousinweb}
\emph{zuìchū} is built from \zh{最} 'most' and \emph{chū} 'the start', both inherited from Middle Chinese, \emph{tswajH} and \emph{tsrhjo} (Guangyun, in Baxter's transcription), which glosses the second 'the beginning': the very start.
\end{cousinweb}

\begin{grammarlens}[title={What you've built}]
One more word for telling what happened.
\end{grammarlens}

\subsection*{Your turn}
\begin{itemize}
\raggedright
  \item \emph{Say it:} \emph{zuìchū}
  \item \emph{Say it:} \emph{zuìchū}, once more
  \item \emph{Say it:} \emph{zuìchū wǒ bù xǐhuan zhèr} — At first I didn't like it here.
  \item \emph{Recall:} say the Mandarin for at long last, then the Mandarin for and so it followed, then say \emph{zuìchū} again
\end{itemize}

\begin{tcolorbox}[breakable,colback=teal!4,colframe=teal!35!black,title={Before you move on}]
What does zuìchū mean? (\textquotedblleft{}At the very start; initially\textquotedblright{}.) Say it once more.
\end{tcolorbox}
\section[\texorpdfstring{\zh{重新} (chóngxīn)}{重新 (chóngxīn)}]{\zh{重新} (chóngxīn) — anew; over again from the start}

\label{lesson:ZH-C295-chongxin}



\begin{quote}
Before the new one: say the Mandarin for and so it followed, then the Mandarin for at the very start.
\end{quote}

\subsection*{You'll want to know: \zh{重新}}
\textbf{\zh{重新}} — \emph{chóngxīn} — \textquotedblleft{}anew; over again from the start\textquotedblright{}.

\begin{quote}
\textbf{\zh{你得重新写}} \emph{nǐ děi chóngxīn xiě} — You'll have to write it over again.
\end{quote}

The new character, in pieces: \zh{新} — thirteen strokes in three pieces: top left, the piece on the left of \zh{站}, five strokes; under it a bar, a vertical with a hook, a short sweep and a dot; on the right two strokes falling to the left, a bar and a long vertical. The next lesson puts it in your hand.

\begin{cousinweb}
\textbf{\zh{重新}} (\emph{chóngxīn}) is built from \zh{重} 'to repeat' and \zh{新} 'new', both inherited from Middle Chinese, \emph{drjowng} and \emph{sin} (Guangyun, in Baxter's transcription), which glosses the first, in this reading, 'to double' and 'to layer'; the same character read \emph{drjowngX}, glossed 'much' and 'thick', is the source of zhòng 'heavy'.
\end{cousinweb}

\begin{grammarlens}[title={What you've built}]
One more word for telling what happened.
\end{grammarlens}

\subsection*{Your turn}
\begin{itemize}
\raggedright
  \item \emph{Say it:} \emph{chóngxīn}
  \item \emph{Say it:} \emph{chóngxīn}, once more
  \item \emph{Say it:} \emph{nǐ děi chóngxīn xiě} — You'll have to write it over again.
  \item \emph{Recall:} say the Mandarin for and so it followed, then the Mandarin for at the very start, then say \emph{chóngxīn} again
\end{itemize}

\begin{tcolorbox}[breakable,colback=teal!4,colframe=teal!35!black,title={Before you move on}]
What does \zh{重新} mean? (\textquotedblleft{}Anew; over again from the start\textquotedblright{}.) Say it once more.
\end{tcolorbox}
\section[\texorpdfstring{\zh{新} (xīn)}{新 (xīn)}]{\zh{新} — two pieces on the left, one on the right}

\label{lesson:ZH-W295-xin-new}



\begin{quote}
The word first, then the character inside it.

\begin{itemize}
\raggedright
  \item \emph{Recall:} say \emph{chóngxīn} — \textbf{R1}, one lesson back, before it decays
\end{itemize}

\begin{quote}
\zh{站} \zh{机}
\end{quote}

\emph{Read it:} each one aloud, with nothing said first — \emph{zhàn}, \emph{jī}
\end{quote}

\subsection*{Script — two pieces on the left, one on the right}
\hlblockfigure{\detokenize{figures/ZH-W295-xin-new-filmstrip.pdf}}{How \zh{新} is written, stroke by stroke}

\begin{quote}
\zh{新}
\end{quote}

\textbf{Three pieces: two stacked on the left, one on the right.} Count the pieces, not the strokes.

\textbf{Top left, five strokes.} The piece on the left of \zh{站}: a dot, a bar, a dot falling to the right and a short stroke falling to the left under it, then a bar across the foot.

\textbf{Bottom left, four strokes.} A bar, a vertical that hooks to the left at the foot, a short sweep to the lower left and a dot to the lower right: the tree piece of \zh{机}, with a hook.

\textbf{The right piece, four strokes.} A short sweep down to the left at the top. A long stroke down from under it that curves to the lower left. A bar across to the right. Then a long vertical down to the foot.

5 + 4 + 4 = thirteen.

\subsection*{Your turn}
\begin{itemize}
\raggedright
  \item \emph{Write it:} \zh{新}, all thirteen, with the model in front of you
  \item \emph{Cover:} write it again from memory, and count the pieces, not the strokes
\end{itemize}

\begin{tcolorbox}[breakable,colback=teal!4,colframe=teal!35!black,title={Before you move on}]
How many pieces in \zh{新}? (\textbf{Three}.) Where have you made the top left piece? (\textbf{On the left of \zh{站}}.) What is the last stroke? (\textbf{A long vertical on the right}.)

Source: \href{https://raw.githubusercontent.com/chanind/hanzi-writer-data/68d10a4b21150cae5e1ebbd223eed289cf32d90c/data/\%E6\%96\%B0.json}{Hanzi Writer Data \zh{新}, pinned stroke paths and medians}.
\end{tcolorbox}
\section[\texorpdfstring{\zh{当天} (dàngtiān)}{当天 (dàngtiān)}]{\zh{当天} (dàngtiān) — that same day; on that very day}

\label{lesson:ZH-C295-dangtian}



\begin{quote}
Before the new one: say the Mandarin for at the very start, then the Mandarin for over again from the start.
\end{quote}

\subsection*{You'll want to know: \zh{当天}}
\textbf{\zh{当天}} — \emph{dàngtiān} — \textquotedblleft{}that same day; on that very day\textquotedblright{}.

\begin{quote}
\textbf{\zh{他当天就回来了}} \emph{tā dàngtiān jiù huílái le} — He came back that same day.
\end{quote}

\begin{cousinweb}
\textbf{\zh{当天}} (\emph{dàngtiān}) is built from \zh{当} and \zh{天} 'a day', both inherited from Middle Chinese, \emph{tangH} and \emph{then} (Guangyun, in Baxter's transcription), which glosses the first, in this departing-tone reading, 'to be in charge'; how it came to mean 'that very (day)' is not settled. Read \emph{tang}, the same character gives dāng.
\end{cousinweb}

\begin{grammarlens}[title={What you've built}]
One more word for telling what happened.
\end{grammarlens}

\subsection*{Your turn}
\begin{itemize}
\raggedright
  \item \emph{Say it:} \emph{dàngtiān}
  \item \emph{Say it:} \emph{dàngtiān}, once more
  \item \emph{Say it:} \emph{tā dàngtiān jiù huílái le} — He came back that same day.
  \item \emph{Recall:} say the Mandarin for at the very start, then the Mandarin for over again from the start, then say \emph{dàngtiān} again
\end{itemize}

\begin{tcolorbox}[breakable,colback=teal!4,colframe=teal!35!black,title={Before you move on}]
What does \zh{当天} mean? (\textquotedblleft{}That same day; on that very day\textquotedblright{}.) Say it once more.
\end{tcolorbox}
\section[\texorpdfstring{\zh{事后} (shìhòu)}{事后 (shìhòu)}]{\zh{事后} (shìhòu) — after the event; in hindsight}

\label{lesson:ZH-C295-shihou}



\begin{quote}
Before the new one: say the Mandarin for over again from the start, then the Mandarin for on that very day.
\end{quote}

\subsection*{You'll want to know: \zh{事后}}
\textbf{\zh{事后}} — \emph{shìhòu} — \textquotedblleft{}after the event; in hindsight\textquotedblright{}.

\begin{quote}
\textbf{\zh{事后我们都笑了}} \emph{shìhòu wǒmen dōu xiào le} — Afterwards we all laughed.
\end{quote}

\begin{cousinweb}
\textbf{\zh{事后}} (\emph{shìhòu}) is built from \zh{事} 'a matter' and \zh{后} 'after', both inherited from Middle Chinese, \emph{dzriH} and \emph{huwX} (Guangyun, in Baxter's transcription), which glosses the second 'first and after'.
\end{cousinweb}

\begin{grammarlens}[title={What you've built}]
One more word for telling what happened.
\end{grammarlens}

\subsection*{Your turn}
\begin{itemize}
\raggedright
  \item \emph{Say it:} \emph{shìhòu}
  \item \emph{Say it:} \emph{shìhòu}, once more
  \item \emph{Say it:} \emph{shìhòu wǒmen dōu xiào le} — Afterwards we all laughed.
  \item \emph{Recall:} say the Mandarin for over again from the start, then the Mandarin for on that very day, then say \emph{shìhòu} again
\end{itemize}

\begin{tcolorbox}[breakable,colback=teal!4,colframe=teal!35!black,title={Before you move on}]
What does \zh{事后} mean? (\textquotedblleft{}After the event; in hindsight\textquotedblright{}.) Say it once more.
\end{tcolorbox}
\section[dāngchǎng]{dāngchǎng — on the spot; there and then}

\label{lesson:ZH-C295-dangchang}



\begin{quote}
Before the new one: say the Mandarin for on that very day, then the Mandarin for after the event.
\end{quote}

\subsection*{You'll want to know: dāngchǎng}
\emph{dāngchǎng} — \textquotedblleft{}on the spot; there and then\textquotedblright{}.

\begin{quote}
\emph{tā dāngchǎng jiù kū le} — He burst into tears there and then.
\end{quote}

Its characters come later; for now it is a word for the ear.

\begin{cousinweb}
\emph{dāngchǎng} is built from \zh{当} 'at' and \emph{chǎng} 'an open ground', both inherited from Middle Chinese, \emph{tang} and \emph{drjang} (Guangyun, in Baxter's transcription), which glosses the first 'to face' and 'to meet' and the second 'a place for offerings' and 'a threshing floor'. Middle Chinese \emph{drjang} would regularly give cháng; the dipping tone of chǎng is irregular.
\end{cousinweb}

\begin{grammarlens}[title={What you've built}]
One more word for telling what happened.
\end{grammarlens}

\subsection*{Your turn}
\begin{itemize}
\raggedright
  \item \emph{Say it:} \emph{dāngchǎng}
  \item \emph{Say it:} \emph{dāngchǎng}, once more
  \item \emph{Say it:} \emph{tā dāngchǎng jiù kū le} — He burst into tears there and then.
  \item \emph{Recall:} say the Mandarin for on that very day, then the Mandarin for after the event, then say \emph{dāngchǎng} again
\end{itemize}

\begin{tcolorbox}[breakable,colback=teal!4,colframe=teal!35!black,title={Before you move on}]
What does dāngchǎng mean? (\textquotedblleft{}On the spot; there and then\textquotedblright{}.) Say it once more.
\end{tcolorbox}
\section[(dialogue)]{Looking back: from In the News to Starting Over}

\label{lesson:ZH-R295-review-290-295}



\begin{quote}
Say the Mandarin for the news on TV and to report on the news.
\end{quote}

\subsection*{Your turn}
\begin{itemize}
\raggedright
  \item \emph{Say it:} the news on TV, to report on the news, a radio station, the presenter of a show, what everyone is talking about
  \item \emph{Say it:} to interview someone for the news, to invent something, an incident in the news, to break out suddenly, to announce officially
  \item \emph{Say it:} to broadcast on the radio, the people watching a show, to broadcast live, the people listening to the radio, to publish an article
  \item \emph{Say it:} a magazine, the daily paper, to make the papers, the front-page story, to follow a story closely
  \item \emph{Say it:} how things stand, to give rise to something, to have an effect on, at long last, and so it followed
  \item \emph{Say it:} at the very start, over again from the start, on that very day, after the event, right there on the spot
  \item \emph{Read it:} \zh{发} \zh{报} \zh{新} — each one aloud, then a word that holds it
\end{itemize}

\begin{tcolorbox}[breakable,colback=teal!4,colframe=teal!35!black,title={Before you move on}]
The news on TV? (\emph{xīnwén}.) To report on the news? (\emph{bàodào}.) A radio station? (\emph{diàntái}.) The presenter of a show? (\emph{zhǔchírén}.) What everyone is talking about? (\textbf{\zh{热点}}, \emph{rèdiǎn}.) To interview someone for the news? (\emph{cǎifǎng}.) To invent something? (\textbf{\zh{发明}}, \emph{fāmíng}.) An incident in the news? (\emph{shìjiàn}.) To break out suddenly? (\textbf{\zh{突发}}, \emph{tūfā}.) To announce officially? (\emph{xuānbù}.) To broadcast on the radio? (\emph{guǎngbō}.) The people watching a show? (\emph{guānzhòng}.) To broadcast live? (\emph{zhíbō}.) The people listening to the radio? (\emph{tīngzhòng}.) To publish an article? (\emph{fābiǎo}.) A magazine? (\emph{zázhì}.) The daily paper? (\textbf{\zh{报}}, \emph{bào}.) To make the papers? (\textbf{\zh{见报}}, \emph{jiànbào}.) The front-page story? (\emph{tóutiáo}.) To follow a story closely? (\emph{guānzhù}.) How things stand? (\emph{qíngkuàng}.) To give rise to something? (\emph{yǐnqǐ}.) To have an effect on? (\emph{yǐngxiǎng}.) At long last? (\emph{zhōngyú}.) And so it followed? (\emph{yúshì}.) At the very start? (\emph{zuìchū}.) Over again from the start? (\textbf{\zh{重新}}, \emph{chóngxīn}.) On that very day? (\textbf{\zh{当天}}, \emph{dàngtiān}.) After the event? (\textbf{\zh{事后}}, \emph{shìhòu}.) Right there on the spot? (\emph{dāngchǎng}.)
\end{tcolorbox}
